\documentclass[12pt, a4paper]{article}

% --- PAQUETES ESENCIALES ---
\usepackage[utf8]{inputenc}
\usepackage[T1]{fontenc}
\usepackage[spanish]{babel}
\usepackage{geometry}
\geometry{top=2.5cm, bottom=2.5cm, left=2.5cm, right=2.5cm}

% --- FUENTE TIMES NEW ROMAN ---
\usepackage{newtxtext,newtxmath}

% --- FORMATO DE PÁRRAFO, LISTAS Y ENLACES ---
\usepackage{parskip} % Espacio limpio entre párrafos sin sangría
\usepackage{ragged2e} % Asegura el justificado completo en todo el documento
\usepackage{enumitem} % Personalización de listas
\usepackage{csquotes}
\usepackage{hyperref}

\hypersetup{
    colorlinks=true,
    linkcolor=blue,
    urlcolor=blue,
    citecolor=blue
}

% --- CONFIGURACIÓN BIBLIOGRÁFICA IEEE ENLAZADA A referencias.bib ---
\usepackage[backend=biber,style=ieee,sorting=none]{biblatex}
\addbibresource{referencias.bib}

% --- INFORMACIÓN DEL DOCUMENTO ---
\begin{document}
\justifying % Aplica alineación justificada a todo el documento

% --- ENCABEZADO EN TAMAÑO 12PT ---
\begin{center}
    \textbf{GA-TL-02 - FORM} \\
    \textbf{Modelado de Comunicación entre Procesos Distribuidos} \\
    \vspace{0.3cm}
\end{center}

\vspace{0.5cm}

% =========================================================================
% PREGUNTA 1
% =========================================================================
\textbf{1. ¿Por qué usan TCP puro sin cifrado para datos bancarios sensibles y cómo evitan que los mensajes se fragmenten o se peguen en la red?}

TCP garantiza la entrega ordenada de un flujo de bytes, pero no cifra la información. Por tanto, utilizar TCP sin TLS en una comunicación que transporta información bancaria sensible no sería una práctica recomendable.

Para proteger esa comunicación se debería implementar TLS, por ejemplo, mediante TLS 1.3, y utilizar autenticación mutua (mTLS) si ambos servicios necesitan verificar sus identidades \cite{vansteen2023}.

Respecto a la fragmentación o unión de mensajes, TCP no trabaja con mensajes individuales, sino con un flujo continuo de bytes \cite{vansteen2023}. Para solucionar esto se necesita un mecanismo de delimitación, como:

\begin{itemize}[leftmargin=1.5cm]
    \item \textbf{Prefijo de longitud:} enviar primero el tamaño del mensaje y después su contenido.
    \item \textbf{Delimitadores:} utilizar un carácter especial que indique dónde termina cada mensaje.
\end{itemize}

De esta manera, el servidor puede reconstruir correctamente cada mensaje recibido.

\vspace{0.4cm}

% =========================================================================
% PREGUNTA 2
% =========================================================================
\textbf{2. Si la transferencia se guarda en PostgreSQL pero RabbitMQ se cae justo antes de publicar el evento, ¿cómo evitan que el cliente se quede sin notificación ni auditoría?}

El problema es que PostgreSQL y RabbitMQ son sistemas independientes. Si la transferencia se confirma en PostgreSQL, pero RabbitMQ falla antes de recibir el evento, la transferencia queda registrada, pero el sistema de notificaciones y auditorías podría no enterarse.

Una solución es implementar el patrón Transactional Outbox \cite{newman2021}.

Este patrón consiste en guardar la transferencia y un evento pendiente de publicación dentro de la misma transacción de PostgreSQL. Después, un proceso independiente consulta los eventos pendientes y los publica en RabbitMQ cuando el servicio esté disponible.

Así se evita perder el evento, aunque RabbitMQ experimente una caída temporal.

\vspace{0.4cm}

% =========================================================================
% PREGUNTA 3
% =========================================================================
\textbf{3. Si el servicio de saldos realiza el débito pero la respuesta gRPC tarda y el orquestador reintenta la petición, ¿cómo garantizan que no se cobre dos veces?}

El problema ocurre cuando el servicio de saldos ejecuta correctamente la operación por gRPC \cite{buzhin2023}, pero la respuesta tarda en llegar. El orquestador puede interpretar que hubo un fallo y volver a enviar la solicitud.

Para evitarlo se utiliza un mecanismo llamado idempotencia \cite{newman2021}.

Consiste en asignar un identificador único a cada transferencia, conocido como idempotency\_key o clave de idempotencia. Cuando el servicio recibe una solicitud:

\begin{enumerate}[leftmargin=1.5cm]
    \item Comprueba si esa clave ya fue procesada.
    \item Si es nueva, ejecuta el débito y registra el resultado.
    \item Si recibe nuevamente la misma clave, no vuelve a ejecutar el débito.
    \item Devuelve el resultado de la operación original.
\end{enumerate}

También se debe garantizar la unicidad de esa clave en la base de datos mediante una restricción UNIQUE y registrar el resultado dentro de la misma transacción.

\vspace{0.4cm}

% =========================================================================
% PREGUNTA 4
% =========================================================================
\textbf{4. Como todo el flujo principal es síncrono, ¿qué pasa si el motor antifraude se vuelve lento o se cae? ¿Se bloquea toda la banca móvil?}

Si el motor antifraude se vuelve lento, aumenta el tiempo de respuesta de las transferencias \cite{blinowski2022}. Si deja de funcionar, las solicitudes que dependen de su validación pueden quedar bloqueadas hasta que se produzca un error o expire el tiempo de espera (timeout).

Esto podría provocar una saturación del servidor y afectar a otras solicitudes de la banca móvil. Para evitarlo se pueden implementar tres mecanismos \cite{newman2021}:

\begin{itemize}[leftmargin=1.5cm]
    \item \textbf{Timeout:} establecer un tiempo máximo de espera para la respuesta del motor antifraude.
    \item \textbf{Circuit Breaker:} si el motor antifraude falla repetidamente, se interrumpen temporalmente las llamadas para evitar sobrecargar el sistema.
    \item \textbf{Bulkhead:} aislar los recursos destinados al motor antifraude para que su caída no afecte a los demás servicios.
\end{itemize}

Además, se debe definir una política de seguridad: si el motor antifraude no está disponible, la transferencia podría rechazarse temporalmente o quedar pendiente de validación, según las reglas del banco.

\vspace{0.4cm}

% =========================================================================
% PREGUNTA 5
% =========================================================================
\textbf{5. Como las flechas van en un solo sentido y la notificación es asíncrona, ¿en qué momento exacto recibe el cliente la confirmación de que su transferencia fue exitosa?}

El servidor de aplicación coordina las validaciones y operaciones necesarias siguiendo esta secuencia:

\begin{enumerate}[leftmargin=1.5cm]
    \item Recibe la solicitud del cliente.
    \item Consulta al motor antifraude.
    \item Ejecuta el débito y crédito mediante gRPC.
    \item Confirma la transacción en PostgreSQL mediante COMMIT.
    \item Publica el evento TransferenciaExitosa en RabbitMQ de manera asíncrona.
\end{enumerate}

El cliente debería recibir la confirmación de éxito cuando el servidor de aplicación haya recibido la confirmación de que PostgreSQL ejecutó correctamente el COMMIT \cite{auer2021}. No necesita esperar a que RabbitMQ publique el evento ni a que el sistema de notificaciones y auditorías lo consuma.

\vspace{0.4cm}

\nocite{*}
\printbibliography[title={Referencias asociadas al diagrama de comunicación}]

\end{document}